\documentclass[11pt,a4paper]{article}
\usepackage[T2A]{fontenc}
\usepackage[utf8]{inputenc}
\usepackage[russian]{babel}
\usepackage[a4paper,left=22mm,right=22mm,top=23mm,bottom=22mm]{geometry}
\usepackage{amsmath,amssymb,mathtools}
\usepackage{array,booktabs,longtable}
\usepackage{xcolor}
\usepackage{tikz}
\usepackage{hyperref}
\hypersetup{colorlinks=true,linkcolor=blue!55!black,urlcolor=blue!55!black}
\setlength{\parindent}{0pt}
\setlength{\parskip}{4pt}
\setlength{\emergencystretch}{2em}
\renewcommand{\arraystretch}{1.3}

\begin{document}

% Титульный лист.
\thispagestyle{empty}
\vspace*{0.27\textheight}
\begin{center}
  {\Huge\bfseries Ревью домашней работы\par}
  \vspace{20mm}
  {\large Автор: Лёвин Артём Александрович\par}
  \vspace{6mm}
  {\large 13 октября 2026 года\par}
\end{center}
\vfill
\begin{center}
\begin{tikzpicture}
  \draw[blue!35, line width=0.85pt]
    (0,0) .. controls (2.1,0.8) and (4.3,-0.8) .. (6.3,0)
          .. controls (8.4,0.8) and (10.6,-0.8) .. (12.7,0);
\end{tikzpicture}
\end{center}
\newpage

\section*{Сводная таблица ошибок}
\small
\setlength{\tabcolsep}{3pt}
\begin{tabular}{@{}p{0.075\linewidth}p{0.16\linewidth}p{0.29\linewidth}p{0.17\linewidth}p{0.17\linewidth}@{}}
\toprule
\textbf{№} & \textbf{Тема} & \textbf{Суть ошибки} & \textbf{Ваш ответ} & \textbf{Правильный ответ} \\
\midrule
\hyperref[task:2]{2} & Доли минуты, перевод в секунды & При вычислении двух пятых минуты число 60 ошибочно разделено на 5 с результатом 2 вместо 12. & $25\,\text{с}>4\,\text{с}$ & $25\,\text{с}>24\,\text{с}$ \\
\bottomrule
\end{tabular}
\normalsize
\clearpage

% Разделитель.
\thispagestyle{empty}
\vspace*{\fill}
\begin{center}
  {\LARGE\bfseries Разбор ошибок}
\end{center}
\vspace*{\fill}
\clearpage

\phantomsection\label{task:2}
\section*{Задание № 2}

\textbf{Текст задания.} Сравните $\frac{5}{12}$ минуты и $\frac25$ минуты, выразив оба промежутка в секундах.

\textbf{Ваше решение.} В записи сначала верно переведены пять двенадцатых минуты в секунды. Затем при переводе двух пятых минуты получено неверное количество секунд:
\[
  \frac{5}{12}\cdot\frac{60}{1}=25\,\text{с},
  \qquad
  \frac25\cdot\frac{60}{1}=4\,\text{с}.
\]
Далее выполнено сравнение:
\[
  25>4,
  \qquad
  \frac{5}{12}>\frac25.
\]

\textbf{Ваш ответ.} $25\,\text{с}>4\,\text{с}$; первая доля минуты больше второй.

\textcolor{red}{\textbf{Ошибка:}} неверное сокращение дроби при переводе двух пятых минуты в секунды. Итоговый знак сравнения выбран верно, однако один из сравниваемых промежутков вычислен неправильно.

\textbf{Где именно возникает ошибка.}

Последний правильный шаг --- запись выражения, которое переводит долю минуты в секунды:
\[
  \frac25\cdot\frac{60}{1}.
\]
Первый неверный шаг --- сокращение чисел 60 и 5: в записи число 60 после сокращения заменено на 2, тогда как должно получиться 12. Из-за этого вместо 24 секунд появились 4 секунды.

\textbf{Почему этот шаг неверен.}

Одна минута содержит 60 секунд. Чтобы найти две пятых минуты, нужно сначала разделить 60 секунд на пять равных частей, а затем взять две такие части. При делении 60 на 5 получается 12, потому что пять раз по 12 дают 60:
\[
  60:5=12,\qquad 12\cdot5=60.
\]
Число 2 не может быть результатом такого деления: пять раз по 2 дают только 10, а не 60:
\[
  2\cdot5=10\ne60.
\]
Следовательно, при сокращении дроби $\frac{60}{5}$ нужно получить 12. Число 2 из числителя исходной дроби $\frac25$ при этом сохраняется: оно показывает, сколько пятых частей минуты необходимо взять. Поэтому правильное вычисление заканчивается умножением 12 на 2, а не 2 на 2.

\textcolor{blue!60!black}{\textbf{Ключевая идея:}} чтобы перевести часть минуты в секунды, умножьте дробь на 60. Для дроби $\frac ab$ это означает: сначала разделить 60 на $b$, если деление выполняется нацело, затем умножить результат на $a$. При сокращении обязательно проверяйте деление обратным умножением.

\textbf{Исправление от последнего правильного шага.}
\[
\begin{aligned}
  \frac25\cdot\frac{60}{1}
  &=\frac{2\cdot60}{5}\\
  &=2\cdot\frac{60}{5}\\
  &=2\cdot12\\
  &=24\,\text{с}.
\end{aligned}
\]
Таким образом, необходимо заменить четыре секунды на двадцать четыре секунды. Ранее правильно найденные 25 секунд сохраняются.

\textbf{Полное правильное решение.}

\textit{Шаг 1.} Переведём первый промежуток в секунды. Знаменатель 12 показывает, на сколько одинаковых частей делим минуту; числитель 5 показывает, сколько таких частей берём:
\[
\begin{aligned}
 \frac{5}{12}\,\text{мин}
 &=\frac{5}{12}\cdot60\,\text{с}\\
 &=5\cdot(60:12)\,\text{с}\\
 &=5\cdot5\,\text{с}\\
 &=25\,\text{с}.
\end{aligned}
\]
Здесь $60:12=5$: одна двенадцатая минуты равна пяти секундам, значит пять двенадцатых равны двадцати пяти секундам.

\textit{Шаг 2.} Переведём второй промежуток в секунды:
\[
\begin{aligned}
 \frac25\,\text{мин}
 &=\frac25\cdot60\,\text{с}\\
 &=2\cdot(60:5)\,\text{с}\\
 &=2\cdot12\,\text{с}\\
 &=24\,\text{с}.
\end{aligned}
\]
Одна пятая минуты равна двенадцати секундам, а две пятых --- двадцати четырём секундам.

\textit{Шаг 3.} Теперь обе величины выражены в одинаковых единицах, поэтому можно непосредственно сравнить числа 25 и 24:
\[
 25\,\text{с}>24\,\text{с}.
\]
Следовательно, пять двенадцатых минуты больше двух пятых минуты:
\[
  \boxed{\frac{5}{12}\,\text{мин}>\frac25\,\text{мин}}.
\]

\textbf{Проверка.} Переведём найденные секунды обратно в минуты:
\[
  \frac{25}{60}=\frac{5}{12},
  \qquad
  \frac{24}{60}=\frac25.
\]
Обе исходные дроби восстановлены правильно. Дополнительно разность промежутков составляет одну секунду:
\[
 25-24=1\,\text{с},
 \qquad
 \frac{5}{12}-\frac25
 =\frac{25}{60}-\frac{24}{60}
 =\frac1{60}\,\text{мин}.
\]
Одна шестидесятая минуты равна одной секунде, поэтому результаты двух способов проверки совпадают.

\textcolor{teal!60!black}{\textbf{Итог:}} знак сравнения найден правильно, но перевод двух пятых минуты в секунды требует исправления. Первый промежуток равен 25 секундам, второй --- 24 секундам. Первый длиннее второго на одну секунду.

\textbf{Задача на закрепление.}
Сравните $\frac{7}{12}$ минуты и $\frac35$ минуты, выразив оба промежутка в секундах.

\vfill
\begin{center}
\footnotesize Лёвин Артём Александрович эксклюзивно для Ксении Васильченко
\end{center}

\end{document}
